\section{Proofs of the consequences}\label{app:consequences}

We first strengthen the count of orders remaining after any comparison
prefix. This gives the production corollary. We then use an ordered
partition supplied by Cardinal et al. to handle initial partial
information, and give an alternative initialization through production.

\subsection{Orders remaining after a comparison prefix}

Rescaling the weights used in the counting experiment of
Section~\ref{sec:movement-volume} leaves the adversary and all refinement
estimates unchanged. Reducing the chance of entering a child
increases the chance of stopping at the required storage node. This
reduces the loss incurred once per item, at an additional cost each time
an item moves. Lemma~\ref{lem:rare-refinements} bounds how often that cost
is incurred.

\begin{lemma}[Orders remaining after any prefix]\label{lem:prefix-orders}
For every sufficiently large fixed $k$, put
$d_k=-\log(1-1/k)=O(1/k)$. There is a nonnegative
$\delta_k=O(\log k/k)$ such that, after any $q$ completed calls in the
recursive phase, the number $N_{\rm ord}$ of represented orders satisfies
\begin{equation}\label{eq:prefix-orders}
\log N_{\rm ord}\ge\log(n!)-nd_k-(1+\delta_k)q.
\end{equation}
The constants are independent of $n$ and of the comparison sequence.
\end{lemma}
\begin{proof}
Use the same experiment as in Lemma~\ref{lem:packing}, with child weights
\[
\widetilde w(c)=\frac2k w(c).
\]
The total child weight is at most $1/k$, so the stopping probability at
each node is at least $1-1/k$. Let $M$ be the number of refinements, and
let $\widetilde{\mathcal W}$ be the product of these new weights along
all item paths. The argument of Lemma~\ref{lem:packing} gives
\[
N_{\rm ord}\ge n!(1-1/k)^n Z\widetilde{\mathcal W},
\qquad
-\log\widetilde{\mathcal W}
=\sum_r\ell(c_r)+M\log(k/2).
\]
Summing the unchanged refinement estimate \eqref{eq:affine-cost} and
using \eqref{eq:total-refinement-gain} yields
\[
\sum_r\ell(c_r)\le a\log Z+aa_\eta q+bM.
\]
Consequently
\[
\log N_{\rm ord}\ge
\log(n!)-nd_k+(1-a)\log Z-aa_\eta q
-\bigl(b+\log(k/2)\bigr)M.
\]
Here $(1-a)\log Z\ge0$, $aa_\eta=1+O(1/k)$,
$b=O(\log k)$, and $M=O(q/k)$ by
Lemma~\ref{lem:rare-refinements}. These bounds prove
\eqref{eq:prefix-orders} with a nonnegative $\delta_k=O(\log k/k)$.
\end{proof}

\subsection{Producing a target order}

\begin{proof}[Proof of Corollary~\ref{cor:production}]
For sufficiently large $n$, let $k=k(n)$ be the parameter selected in
Section~\ref{sec:parameter-selection}. If the production algorithm reaches
the horizon $Q(n)$, then it has made at least $n\log n\ge I(R)$ comparisons.
Otherwise it finishes after $q<Q(n)$ completed calls. Its output assigns
each key to one slot of $R$. Correctness requires that every total order
consistent with its comparison transcript respect this assignment.
Exactly $e(R)$ total orders respect the assignment, and all represented
orders are consistent with the transcript. Thus $N_{\rm ord}\le e(R)$, and
Lemma~\ref{lem:prefix-orders} gives
\begin{equation}\label{eq:production-bound}
q\ge\frac{I(R)-nd_{k(n)}}{1+\delta_{k(n)}}.
\end{equation}
Both $d_{k(n)}$ and $\delta_{k(n)}$ tend to zero. Taking
$\varepsilon_n=\max\{d_{k(n)},\delta_{k(n)}\}$ and using
$(1+\delta)^{-1}\ge1-\delta$ gives the stated bound. Set
$\varepsilon_n=1$ for the finitely many smaller input sizes, making their
bound nonpositive. These choices are independent of $R$. The adversary's
legality and running time are those of Theorem~\ref{thm:main}.
\end{proof}

\subsection{A compatible partition into layers}

The reduction for an initial partial order needs an ordered partition
whose layers retain many permutations. We use the following consequence
of Cardinal et al.'s construction.

\begin{lemma}[Weak-order extension]\label{lem:weak-order-extension}
Given a partial order $P$ on $n$ elements, one can deterministically
construct in polynomial time a weak order $W$ extending $P$ such that
\[
\log e(W)\ge\log e(P)-C_0 n\log\log(n+2)
\]
for an absolute constant $C_0$.
\end{lemma}
\begin{proof}
Cardinal et al.~\cite[Section~3, Lemmas~1 and~3]{cardinal-production}
construct $W$ in polynomial time with
\[
I(W)\le I(P)+2n\log\bigl(I(P)/n+1\bigr)+O(n).
\]
Since $I(P)\le\log(n!)\le n\log n$, the extra term is
$O(n\log\log(n+2))$. Substitute
$I(W)-I(P)=\log e(P)-\log e(W)$.
\end{proof}

An arbitrary partition into antichains need not admit a compatible layer
order. For example, the relations $a<b$ and $c<d$ permit the antichains
$\{a,d\}$ and $\{b,c\}$, but require opposite orders between them.
Lemma~\ref{lem:weak-order-extension} supplies both the partition and a
layer order respecting all initial relations.

\begin{proof}[Proof of Corollary~\ref{cor:initial-poset}]
Construct $W$ using Lemma~\ref{lem:weak-order-extension}, and let its
layers have sizes $m_1,\ldots,m_t$. Choose $k=k(n)$ once, using the full
input size. Within every layer of size at least two, run a copy of the
ordinary adversary with these parameters and its local horizon
$Q(m_i)=m_i\ceil{\log m_i}$. A singleton needs no comparisons. Answer
comparisons between layers from their fixed order.

Each local transcript has a consistent total order. Concatenating these
orders gives a linear extension of $W$ and hence of $P$, proving legality.
At termination, each layer must have a unique possible order: two possible
orders in one layer, combined with the same orders in the others, would
give two global orders consistent with the full transcript.

Let $q_i$ count comparisons within layer $i$. If its horizon is reached,
$q_i\ge m_i\log m_i\ge\log(m_i!)$. Otherwise
Lemma~\ref{lem:prefix-orders}, with $N_{\rm ord}=1$, gives
\[
q_i\ge\frac{\log(m_i!)-m_i d_k}{1+\delta_k}.
\]
This inequality also holds for a singleton and for a layer whose horizon
was reached. Since $\sum_i m_i=n$ and $e(W)=\prod_i m_i!$, summing gives
\begin{align*}
q\ge\sum_i q_i
&\ge\frac{\log e(W)-nd_k}{1+\delta_k}\\
&\ge\frac{\log e(P)-C_0 n\log\log(n+2)-nd_k}{1+\delta_k}.
\end{align*}
As $k(n)\to\infty$, this implies the stated bound, uniformly over $P$.

The common choice of $k(n)$ is essential to this estimate: each layer
has the same $d_k$ and $\delta_k$, even when its size stays small.
It also preserves polynomial time in $n$. Each finite calculation
enumerates at most $E_k\le n$ assignments. Across all layers, the horizons
sum to at most $n\ceil{\log n}$, and the refinement and precision bounds
of Section~\ref{sec:runtime} apply with the ambient size $n$.
The construction of $W$ is polynomial-time as well. If the parameter
selection produces no eligible tuple, use a fixed topological order of
$P$; enlarging the absolute error constant covers these finitely many
input sizes.
\end{proof}

\subsection{Initialization by partial-order production}\label{sec:production-initialization}

There is also a direct way to initialize from $P$, using a production
algorithm as a black box. It gives the bound
\[
q\ge\log e(P)-r(n),\qquad r(n)=o(n\log n),
\]
with one function $r$ for all $n$-element initial orders.

Run Cardinal et al.'s deterministic production algorithm on target $P$,
answering its comparisons with our ordinary adversary on temporary keys.
Its comparison phase uses the deterministic multiple-selection algorithm
of Kaligosi et al.~\cite{multiple-selection}.
The producer returns an assignment of those keys to the elements of $P$.
Save the adversary's state. Translate each subsequent query through this
assignment and continue the same adversary. Every order consistent with
the initialization transcript respects the returned assignment, so all
later answers are legal for $P$.

The producer followed by any correct sorter for $P$ is a correct ordinary
sorting algorithm. Write the ordinary adversary's guarantee as
$\log(n!)-s(n)$, where $s(n)=o(n\log n)$. The production upper bound
\cite[Theorem~4]{cardinal-production} gives
$U(P)\le I(P)+u(n)$ with $u(n)=o(n\log n)$ independent of $P$:
for each fixed $\delta>0$ its bound is at most
$(1+\delta)I(P)+C_\delta n$, and $I(P)\le n\log n$.
Subtracting the initialization comparisons gives
\[
q\ge\log(n!)-s(n)-U(P)
\ge\log e(P)-s(n)-u(n).
\]
Both the initialization and every subsequent response take polynomial
time. This construction needs no computation of $e(P)$.
